\section{Conclusion}
Incorporating commonsense knowledge into models is becoming a popular trend as it is important for our models to mimic humans and the way they utilize commonsense knowledge in performing different tasks. One danger of mimicking humans is adopting their biases. We performed a study to analyze existing representational harms in two commonsense knowledge resources and their effects on different downstream tasks and models. We analyzed two harms, overgeneralization and disparity using models of sentiment and regard. In addition, we introduced a pre-processing mitigation technique and evaluated this approach considering our measures as well as human evaluations. Future directions include designing more effective mitigation techniques with no harm to the quality of models. 

\section*{Acknowledgments}
We thank anonymous reviewers for providing insightful feedback along with Brendan Kennedy and Lee Kezar for their comments and help. Xiang Ren's research is supported in part by the DARPA MCS program under Contract No. N660011924033, the Defense Advanced Research Projects Agency with award W911NF-19-20271, NSF IIS 2048211, and NSF SMA 182926. This material is based upon work supported, in part, by the Defense Advanced Research Projects Agency (DARPA) and Army Research Office (ARO) under Contract No. W911NF-21-C-0002.

\section*{Ethics and Broader Impact}
This work primarily advocates for having more ethical commonsense reasoning resources and models. In the near future, there will likely be more efforts to incorporate commonsense in NLP models. Conflating human biases with commonsense is harmful. Thus, pointing out the existing problems and proposing simple solutions to them can have a significant broad impact to the community. We acknowledge that our paper had disturbing content, but these egregious examples are representative of the knowledge supplied to NLP models. Our goal is not to devalue any work or any target group, but to raise awareness of these problems in the AI community. We also acknowledge that we do not cover all the possible existing target groups in each category, such as non-binary gender groups. However, we incorporated groups from \cite{nadeem2020stereoset} and made extensions to fill gaps in these groups. Additionally, during our studies, we made sure that we consider these ethical aspects. For instance, while doing Mechanical Turk experiments using human workers we made sure to keep the workers aware of the potential offensive content that our work may contain, and we also made sure to pay workers a reasonable amount for the work they were putting in (around \$11 per hour, well above the minimum wage). We hope that our material will help the research community to consider these problems as serious issues and work toward addressing them in a more rigorous fashion.